\documentclass[11pt]{article}
\pagestyle{plain}

% \pgfplotsset{compat=1.18}
\usepackage[margin=1in]{geometry}
\usepackage{amsmath,amsfonts,amssymb,amsthm}
\usepackage{graphicx}
\usepackage{enumerate}
\usepackage{bbm}
\usepackage{verbatim}
\usepackage{hyperref,color}
\usepackage[capitalize,nameinlink]{cleveref}
\usepackage[dvipsnames]{xcolor}
\hypersetup{
	colorlinks=true,
	pdfpagemode=UseNone,
	citecolor=OliveGreen,
	linkcolor=NavyBlue,
	urlcolor=Magenta,
	pdfstartview=FitW
}
\usepackage{appendix}
\usepackage{makecell}
\usepackage{tablefootnote}
\crefname{appsec}{Appendix}{Appendices}
\usepackage{tikz}
\usepackage{pgfplots}
\usepackage{xifthen}
\usepackage{subcaption}
\usepackage{hhline}

\theoremstyle{plain}
\newtheorem{theorem}{Theorem}[section]
\newtheorem{proposition}[theorem]{Proposition}
\newtheorem{lemma}[theorem]{Lemma}
\newtheorem{corollary}[theorem]{Corollary}
\newtheorem{conjecture}[theorem]{Conjecture}
\newtheorem{problem}[theorem]{Problem}
\newtheorem{claim}[theorem]{Claim}
\newtheorem{fact}[theorem]{Fact}

\theoremstyle{definition}
\newtheorem{definition}[theorem]{Definition}
\newtheorem{example}[theorem]{Example}
\newtheorem{observation}[theorem]{Observation}
\newtheorem{assumption}[theorem]{Assumption}
\newtheorem*{assumption*}{Assumption}

\theoremstyle{remark}
\newtheorem{remark}[theorem]{Remark}

\crefname{lemma}{Lemma}{Lemmas}
\crefname{theorem}{Theorem}{Theorems}
\crefname{definition}{Definition}{Definitions}
\crefname{fact}{Fact}{Facts}
\crefname{claim}{Claim}{Claims}
\crefname{proposition}{Proposition}{Propositions}

\usetikzlibrary{arrows.meta,positioning}

\newcommand{\dif}{\,\mathrm{d}}
\newcommand{\Cov}{\mathrm{Cov}}
\newcommand{\MI}{\mathrm{I}}
\newcommand{\ind}{\*1}
\newcommand{\allone}{\mathbf{1}}
\newcommand{\tr}{\mathrm{tr}}
\newcommand{\wt}{\mathrm{weight}}
\DeclareMathOperator*{\argmax}{arg\,max}
\newcommand{\norm}[1]{\left\lVert #1 \right\rVert}
\newcommand{\ceil}[1]{\left\lceil #1 \right\rceil}
\newcommand{\floor}[1]{\left\lfloor #1 \right\rfloor}
\newcommand{\ip}[2]{\left\langle #1 , #2 \right\rangle}
\newcommand{\grad}{\nabla}
\newcommand{\hessian}{\nabla^2}
\newcommand{\trans}{\intercal}
\newcommand{\diag}{\mathrm{diag}}
\newcommand{\poly}{\mathrm{poly}}
\newcommand{\dist}{\mathrm{dist}}
\newcommand{\sgn}{\mathrm{sgn}}
\newcommand{\fpras}{\mathsf{FPRAS}}
\newcommand{\fpaus}{\mathsf{FPAUS}}
\newcommand{\fptas}{\mathsf{FPTAS}}

\newcommand{\eps}{\varepsilon}

\newcommand{\N}{\mathbb{N}}
\newcommand{\Z}{\mathbb{Z}}
\newcommand{\Q}{\mathbb{Q}}
\newcommand{\R}{\mathbb{R}}
\newcommand{\C}{\mathbb{C}}

\renewcommand{\AA}{\mathcal{A}}
\newcommand{\BB}{\mathcal{B}}
\newcommand{\CC}{\mathcal{C}}
\newcommand{\DD}{\mathcal{D}}
\newcommand{\EE}{\mathcal{E}}
\newcommand{\FF}{\mathcal{F}}
\newcommand{\GG}{\mathcal{G}}
\newcommand{\HH}{\mathcal{H}}
\newcommand{\II}{\mathcal{I}}
\newcommand{\JJ}{\mathcal{J}}
\newcommand{\KK}{\mathcal{K}}
\newcommand{\LL}{\mathcal{L}}
\newcommand{\MM}{\mathcal{M}}
\newcommand{\NN}{\mathcal{N}}
\newcommand{\OO}{\mathcal{O}}
\newcommand{\PP}{\mathcal{P}}
\newcommand{\QQ}{\mathcal{Q}}
\newcommand{\RR}{\mathcal{R}}
\renewcommand{\SS}{\mathcal{S}}
\newcommand{\TT}{\mathcal{T}}
\newcommand{\UU}{\mathcal{U}}
\newcommand{\VV}{\mathcal{V}}
\newcommand{\WW}{\mathcal{W}}
\newcommand{\XX}{\mathcal{X}}
\newcommand{\YY}{\mathcal{Y}}
\newcommand{\ZZ}{\mathcal{Z}}

\newcommand{\KL}[2]{D_{\mathrm{KL}} \left( #1 \,\Vert\, #2 \right)}
\newcommand{\KLbig}[2]{D_{\mathrm{KL}} \left( #1 \,\big\Vert\, #2 \right)}
\newcommand{\KLBig}[2]{D_{\mathrm{KL}} \left( #1 \,\Big\Vert\, #2 \right)}

\newcommand{\SKL}[2]{D_{\mathrm{SKL}} \left( #1 \,\Vert\, #2 \right)}
\newcommand{\SKLbig}[2]{D_{\mathrm{SKL}} \left( #1 \,\big\Vert\, #2 \right)}
\newcommand{\SKLBig}[2]{D_{\mathrm{SKL}} \left( #1 \,\Big\Vert\, #2 \right)}

\newcommand{\qandq}{\quad\text{and}\quad}
\newcommand{\supp}{\mathrm{supp}}

\newcommand{\DKL}[2]{\-D_{\-{KL}}\left(#1 \parallel #2\right)}
\newcommand{\TV}[2]{d_{\mathrm{TV}}\left({#1},\,{#2}\right)}
\newcommand{\e}{\mathrm{e}}
\renewcommand{\epsilon}{\varepsilon}
\renewcommand{\emptyset}{\varnothing}
\newcommand{\set}[1]{\left\{#1\right\}}
\newcommand{\tuple}[1]{\left(#1\right)} 
\newcommand{\inner}[2]{\left\langle #1,#2\right\rangle}
\newcommand{\tp}{\tuple}
\newcommand{\ol}{\overline}
\newcommand{\abs}[1]{\left\vert#1\right\vert}
\newcommand{\ctp}[1]{\left\lceil#1\right\rceil}
\newcommand{\ftp}[1]{\left\lfloor#1\right\rfloor}
\newcommand{\Ind}[1]{\-{Ind}_{#1}}
\newcommand{\todo}[1]{{\color{red} TODO: #1}}

\newcommand{\Cor}{\Psi^{\mathrm{sym}}}
\newcommand{\cor}{\Psi^{\-{cor}}}
%\newcommand{\Inf}{\Psi}

\def\*#1{\boldsymbol{#1}} % Use \*A for \mathbf{A}
\def\+#1{\mathcal{#1}} % Use \+A for \mathcal{A}
\def\-#1{\mathrm{#1}} % Use \-A for \mathrm{A}
\def\=#1{\mathbb{#1}} % Use \^A for \mathbb{A}
\def\!#1{\mathfrak{#1}} % Use \!A for \mathfrak{A}

\newcommand*\diff{\mathop{}\!\mathrm{d}}

% probability
% \DeclareMathOperator*{\oPr}{\mathbf{Pr}}
%\def\oPr{\mathop{\mathrm{Pr}}}
\def\oPr{\mathrm{Pr}}
\renewcommand{\Pr}[2][]{ \ifthenelse{\isempty{#1}}
  {\oPr\left[#2\right]}
  {\oPr_{#1}\left[#2\right]} } % Use \Pr[a]{b} for \mathbf{Pr}_a[b], \Pr{b} for  \mathbf{Pr}[b]

% expectation: relies on the package "dsfont"
% \DeclareMathOperator*{\oE}{\mathbb{E}}
%\def\oE{\mathop{\mathbb{E}}}
\def\oE{\mathbb{E}}
\newcommand{\E}[2][]{ \ifthenelse{\isempty{#1}}
  {\oE\left[#2\right]}
  {\oE_{#1}\left[#2\right]} }

% variance
%\DeclareMathOperator*{\oVar}{\mathbf{Var}}
\def\oVar{\mathrm{Var}}
\newcommand{\Var}[2][]{ \ifthenelse{\isempty{#1}}
  {\oVar\left(#2\right)}
  {\oVar_{#1}\left(#2\right)} }

% entropy
% \DeclareMathOperator*{\oEnt}{\mathbf{Ent}}
\def\oEnt{\mathrm{Ent}}
\newcommand{\Ent}[2][]{ \ifthenelse{\isempty{#1}}
  {\oEnt\left[#2\right]}
  {\oEnt_{#1}\left[#2\right]} }

\newcommand{\PhiEnt}[2][]{ \ifthenelse{\isempty{#1}}
  {\oEnt^\phi\left[#2\right]}
  {\oEnt^\phi_{#1}\left[#2\right]} }


\newenvironment{Exampleparagraph}[1][]
{
    \par 
    \noindent 
    \textbf{\color{blue} #1}
    \par 
    \vspace{0.5em} 
    \itshape 
}
{
    \par 
    \vspace{1em} 
}

\newcommand{\mathsc}[1]{{\normalfont\textsc{#1}}}

\newcommand{\Sch}{\mathrm{Sch}}
\newcommand{\LSI}{\mathsc{lsi}}
\newcommand{\MLSI}{\mathsc{mlsi}}
\newcommand{\BE}{\mathsc{be}}
\newcommand{\RBE}{\mathrm{r}\mathsc{be}}
\newcommand{\WGE}{\mathsc{wge}}
\newcommand{\SGE}{\mathsc{sge}}
\newcommand{\CSGE}{\mathsc{csge}}

\newcommand{\Inf}{\mathsf{Inf}}
\newcommand{\SurfArea}{\mathsf{SA}}
\newcommand{\SqInf}{\mathsf{SqInf}}
\newcommand{\cube}{\{\pm 1\}^n}
\newcommand{\unif}{\nu_n}
\newcommand{\placeholder}[1]{%
  \begin{center}
    \fbox{\parbox{0.90\textwidth}{\strut\textbf{#1}\strut}}
  \end{center}
}



\newcommand{\ZJ}[1]{{\color{red}Zejia: #1}}
\newcommand{\ZC}[1]{{\color{violet}Zongchen: #1}}
\newcommand{\XZ}[1]{{\color{blue}Xinyuan: #1}}

\title{Witten Laplacian}
\author{}
\date{\today}

\pgfplotsset{compat=1.18}

\begin{document}

\maketitle

\tableofcontents

\newpage

\section{Setup}

\(L = \sum_{i=1}^n \delta_i\) where \(\delta_i g(x) = q_i(x) (g(x^i) - g(x))\).
Also, assume \(q_i(x) = \sqrt{\mu(x^i)}/\sqrt{\mu(x)}\).

\section{Witten Laplacian (discrete version)}

For \(f\) with mean \(0\), let \(g = (-L)^{-1} f\). Then
\[
Var_\mu(f) = E_\mu(f (-L) g) = E_\mu(\langle \nabla f,\nabla g\rangle_q).
\]
A straightforward calculation yields
\[
\delta_i f = \delta_i (-L) g = -\sum_{j} \delta_i \delta_j g = - L \delta_i g + \sum_{j} [\delta_j,\delta_i] g.
\]
Note that
\([\delta_j,\delta_i] g(x) = (q_j(x^i) - q_j(x)) \delta_i g - (q_i(x^j) - q_i(x)) \delta_j g\).

Therefore, \(\nabla f\) can be expressed by
\[
\nabla f = (-L \otimes I_n + \mathcal{M}) \nabla g,
\]
where \(\mathcal{M}(x) = \mathrm{diag}(\sum_j (q_j(x^i) - q_j(x))) - (q_i(x^j)-q_i(x))_{1 \le i,j \le n}\).

Therefore,
\[
E_\mu(\langle \nabla f, \nabla g\rangle_q) = E_\mu(\langle \nabla f, (-L \otimes I_n + \mathcal{M})^{-1} \nabla f \rangle_q).
\]

For product distribution on Boolean hypercube, we have
\[
E_\mu(\langle \nabla f, (-L \otimes I_n + \mathcal{M})^{-1} \nabla f \rangle_q) = E_\mu(\langle \nabla f, (-L \otimes I_n)^{-1} \nabla f \rangle_q) = \frac{1}{2} \int_0^{+\infty} \sum_i E_\mu((\partial_i f)(P_t \partial_i f)) dt.
\]

\section{Comment}
 
It seems we cannot expect \(\mathcal{M} \ge \rho I\) as in the Langevin dynamics.

\ZC{(old) My conversation with Gemini says that $-L \otimes I_n + \mathcal{M} \succeq \rho \mathrm{id}$ does not hold pointwisely (i.e., ${\langle \nabla f(x), (-L \otimes I_n + \mathcal{M}) \nabla f(x) \rangle_q} \ge \rho {\langle \nabla f(x), \nabla f(x) \rangle_q}$ does not hold for all $x$), but holds in expectation (i.e., $\E[]{\langle \nabla f, (-L \otimes I_n + \mathcal{M}) \nabla f \rangle_q} \ge \rho \E[]{\langle \nabla f, \nabla f \rangle_q}$).}

\ZC{(old) Want to find $\mathcal{M'}$ such that
\begin{align*}
    \E[]{\langle \nabla f, (-L \otimes I_n + \mathcal{M}) \nabla f \rangle_q} \ge \E[]{\langle \nabla f, \mathcal{M'}\nabla f \rangle_q}
\end{align*}
and then, under similar Dobrushin-type conditions, it holds pointwisely
\begin{align*}
    \mathcal{M'}(x) \succeq \rho I.
\end{align*}
This gives a Brascamp--Lieb inequality with $\mathcal{M'}$.
}

\bigskip
\ZC{(new) Define the \textit{non-directional} derivatives and gradient as
\begin{align*}
    \overline{\grad} f = 
    \begin{pmatrix}
        \overline{\delta}_1 f \\
        \vdots \\
        \overline{\delta}_n f
    \end{pmatrix},
\end{align*}
where
\begin{align*}
    \overline{\delta}_i f(x)
    = q_i(x) \frac{f(x^i)-f(x)}{-x_i} 
    = q_i(x) (-x_i) (f(x^i)-f(x)).
\end{align*}
Then we have
\begin{align*}
    \overline{\grad} Lf = (L \otimes I_n - \mathcal{H}) \overline{\grad} f,
\end{align*}
where $\mathcal{H}$ is defined by a collection of local Hessian matrices $\{H(x) \in \R^{n \times n}: x \in \{-1,+1\}^n\}$ such that
\begin{align*}
    \mathcal{H} v(x) = H(x) v(x),
\end{align*}
for any (gradient) vector field $v: \{-1,+1\}^n \to \R^n$.
If $q_i(x) = \frac{1}{2}$ at all $x$ and $i$, then $H(x) = I$ at all $x$ and $\mathcal{H} = \mathrm{id}$.
}

\ZC{$$\mathbb{E}_\mu [\langle \grad f, (-L\otimes I_n) \grad f \rangle_q] \ge 0$$}

\ZC{$$\langle v, \mathcal{A} v \rangle \ge \langle v, \mathcal{B} v \rangle > 0, \forall v
\quad \Longrightarrow \quad
\langle v, \mathcal{A}^{-1} v \rangle \le \langle v, \mathcal{B}^{-1} v \rangle, \forall v
$$
where $\langle u, v \rangle = \mathbb{E}_\mu [\langle u,v \rangle_q]$.\\
$\mathcal{A} \succeq \mathcal{B} \succ 0
\quad \not\Longrightarrow \quad 
\mathcal{A}^2 \succeq \mathcal{B}^2$
}


\end{document}
